\documentclass[12pt]{book}

\usepackage[utf8]{inputenc}
\usepackage[T1]{fontenc}
\usepackage{lmodern}
% setta i margini per tener conto della rilegatura 
\usepackage[
    a4paper,
    margin=3cm,
    bindingoffset=1cm,
    includeheadfoot
]{geometry}
\usepackage{graphicx}

\usepackage{setspace}
\usepackage[per-mode=symbol]{siunitx}
\usepackage{url}
\usepackage{xurl}
\usepackage{csquotes}
\DeclareQuoteAlias[quotes]{italian}{italian}
\sisetup{list-final-separator = {, and }}
\usepackage[main=english,italian]{babel}
\usepackage[capitalise]{cleveref}
\crefname{section}{Sect.}{Sect.}
\crefname{figure}{Fig.}{Fig.}
\crefname{table}{Tab.}{Tab.}
\crefname{equation}{Eq.}{Eq.}
\Crefname{table}{Table}{Tables}
\Crefname{equation}{Equation}{Equations}
\crefname{equation}{Equation}{Equations}
\def\figname{\csname cref@figure@name\endcsname\xspace}
\def\tabname{\csname cref@table@name\endcsname\xspace}
\def\secname{\csname cref@section@name\endcsname\xspace}
\def\eqname{\csname cref@equation@name\endcsname\xspace}
\def\eqpname{\csname cref@equation@name@plural\endcsname\xspace}
\usepackage{xspace}
\usepackage{hyphenat}
\usepackage{listings}
\usepackage{bytefield}

% setup dei colori, nella libreria xcolor ce ne sono tantissimi altri
\usepackage{color,soul}
\usepackage{xcolor} 
\definecolor{light-gray}{gray}{0.97} 
\definecolor{codegreen}{rgb}{0,0.6,0}

% permette di usare il comando \todo{testo} per scrivere del testo in una casella arancione per indicare i to-do
\usepackage{todonotes}
\presetkeys{todonotes}{inline}{}

% con il comando \lstinputlistings{file.txt} si può inserire codice nel testo
\lstset{
language=Matlab,
breaklines=true,
showstringspaces=false,
backgroundcolor = \color{light-gray},
basicstyle=\fontsize{8}{8}\selectfont\ttfamily,
tabsize=4,
numbers=left,
commentstyle = \color{blue},
keywordstyle = \color{red},
stringstyle = \color{codegreen},
rulecolor = \color{black},
captionpos=b
}

\usepackage[backend=biber,style=ieee,doi=false,isbn=false,mincitenames=1,maxcitenames=4]{biblatex}
\DeclareFieldFormat{sentencecase}{#1} % never apply sentence casing, even if bibtex field is unprotected
\DeclareFieldFormat{titlecase}{#1} % never apply title casing, even if bibtex field is unprotected
\addbibresource{bibliography.bib}
\DeclareFieldFormat{titlecase}{#1} % never apply title casing, even if bibtex field is unprotected
\addbibresource{bibliography.bib}
\usepackage{xpatch}
\xpatchbibmacro{textcite}{\addspace}{\addnbspace}{}{}
\xpatchbibmacro{Textcite}{\addspace}{\addnbspace}{}{}
\setlength\biblabelsep{.6em}
\renewcommand{\bibfont}{\footnotesize}
\DefineBibliographyStrings{english}{
	andothers = et~al\adddot\addspace
}

% indica al comando \includegraphics dove andare a prendere le immagini, ovviamente si può cambiare il path a piacere se non si vuole usare ./images/
\graphicspath{{./images/}}

\usepackage{setspace}

% le prossime sezioni (escluso DOCUMENT) possono essere inseriti o tolti a seconda di come piace di più esteticamente la tesi (o a seconda delle indicazioni del relatore)

%%%%%%%%%%%%%%%%%%%%%%%%%%%CHAPTER%%%%%%%%%%%%%%%%%%%%%%%%%%%
\usepackage{titlesec}
\titleformat{\chapter}{\normalfont\fontsize{18}{27}\bfseries}{\thechapter}{1.5em}{}
\titlespacing*{\chapter}{0pt}{5.5ex plus 1ex minus .2ex}{5.3ex plus .2ex}

%%%%%%%%%%%%%%%%%%%%%%%%%%%SECTION%%%%%%%%%%%%%%%%%%%%%%%%%%%
\titleformat{\section}{\normalfont\fontsize{16}{24}\bfseries}{\thesection}{1em}{}
\titlespacing*{\section}{0pt}{2.5ex plus 1ex minus .2ex}{2.3ex plus .2ex}

%%%%%%%%%%%%%%%%%%%%%%%%%%%SUBSECTION%%%%%%%%%%%%%%%%%%%%%%%%%%%
\titleformat{\subsection}{\normalfont\fontsize{13}{19}\bfseries}{\thesubsection}{1em}{}
\titlespacing*{\subsection}{0pt}{1.5ex plus 1ex minus .2ex}{1.3ex plus .2ex}

%%%%%%%%%%%%%%%%%%%%%%%%%%%SUBSUBSECTION%%%%%%%%%%%%%%%%%%%%%%%%%%%
\titleformat{\subsubsection}{\normalfont\fontsize{12}{18}\itshape\bfseries}{\thesubsubsection}{1em}{}
\titlespacing*{\subsection}{0pt}{1.5ex plus 1ex minus .2ex}{1.3ex plus .2ex}

%%%%%%%%%%%%%%%%%%%%%%%%%%%HEADER%%%%%%%%%%%%%%%%%%%%%%%%%%%
\usepackage{fancyhdr}

\fancypagestyle{plain}{%
    \fancyhf{}% Clear header and footer
    \fancyfoot[C]{
        \if\relax\thepage\else%
            --- {\thepage} ---
        \fi
    }
    \renewcommand{\headrulewidth}{0pt}
    \renewcommand{\footrulewidth}{0pt}
}
\pagestyle{plain}

%%%%%%%%%%%%%%%%%%%%%%%%%%%%%%%%%%%%%%%%%%%%%%%%%%%%%%%%%%%%%
%                       DOCUMENT                            %
%%%%%%%%%%%%%%%%%%%%%%%%%%%%%%%%%%%%%%%%%%%%%%%%%%%%%%%%%%%%%
\begin{document}
    \pagenumbering{gobble}
    \begin{otherlanguage}{italian}
        \input{frontespice}
    \end{otherlanguage}
    \cleardoublepage
    \pagenumbering{Roman}
    \setstretch{1.5}
    
    \begin{otherlanguage}{italian}
        \chapter*{Sommario}
Questa tesi affronta il problema dell'affidabilità delle comunicazioni Vehicle-to-Everything (V2X) in scenari di guida cooperativa, con particolare riferimento al platooning. In tali contesti, la qualità del canale radio e la variabilità delle condizioni di propagazione possono compromettere la trasmissione tempestiva delle informazioni di controllo, con impatti diretti su stabilità, comfort e sicurezza del convoglio. L'obiettivo del lavoro è progettare e valutare un approccio multi-canale che integri più tecnologie di comunicazione, selezionando dinamicamente il canale più affidabile in funzione delle condizioni osservate.

La parte sperimentale unisce la simulazione a eventi discreti dei framework OMNeT++, Veins e Plexe con la simulazione a tempo continuo di Sumo. In primo luogo, viene descritto il processo di preparazione e configurazione dell'ambiente di simulazione, includendo la riproducibilità degli esperimenti e la gestione delle dipendenze software. Successivamente, viene importata e adattata nel nuovo ambiente la logica di Packet Delivery Ratio (PDR) e di SafeSwitch, originariamente sviluppata nel progetto Segata2021Plexe, insieme agli scenari sperimentali di riferimento.

Un contributo centrale del lavoro consiste nella migrazione del modulo plexe\_hetnet dall'utilizzo di SimuLTE a Simu5G, al fine di supportare tecnologie cellulari di nuova generazione e abilitare valutazioni più aderenti agli standard recenti. Infine, vengono eseguite campagne di simulazione per diversi controller di platooning e parametri di traffico, e vengono generati grafici e metriche quantitative (ad esempio, PDR e indicatori legati alla dinamica del convoglio) per confrontare le prestazioni delle diverse configurazioni e valutare l'efficacia delle politiche di selezione del canale.

I risultati mostrano come l'adozione di un approccio multi-canale con meccanismi di stima della qualità del link e commutazione controllata possa migliorare la robustezza della comunicazione in condizioni degradate, permettendo ai veicoli del platoon di mantenere una minore distanza tra loro. La tesi si conclude con un'analisi critica dei limiti del modello e con possibili estensioni future, tra cui l'introduzione di ulteriori canali, strategie di decisione più sofisticate e una validazione su scenari più complessi.
    \end{otherlanguage}
    
    \cleardoublepage
    \chapter*{Summary}
This thesis studies the reliability of V2X communications in platooning scenarios, with the goal of understanding how to safely break up a platoon when link quality degrades. It builds on \emph{SafeSwitch}, a mechanism that links communication quality, controller selection, and inter-vehicle distance in order to avoid abrupt and potentially dangerous transitions.

The main contribution is the reconstruction of a multi-technology experimental environment built around OMNeT++, Veins, Plexe, SUMO, and Simu5G. Within this framework, the scenarios and metrics of the reference work were reproduced, the PDR monitoring logic and the \emph{SafeSwitch} state machine were reimplemented, and, most importantly, the cellular component was migrated from SimuLTE to Simu5G, updating the system to a more current 5G model.

The results show that the mechanism behaves consistently even on the new stack: in critical scenarios \emph{SafeSwitch} performs better than a \emph{naive} fallback, avoiding overly abrupt transitions and maintaining control of the platoon, while in the realistic scenario the migration to 5G introduces no functional regressions and, if anything, leads to fewer losses during cell handovers. The thesis therefore confirms that the architecture can be ported to a more modern toolchain and highlights how communication and control must be designed jointly to guarantee system safety.


    \cleardoublepage
    \tableofcontents
    \cleardoublepage
    
    \pagenumbering{arabic}

    \chapter{Prova}

Qui si inserisce il contenuto; è comodo mettere
\begin{verbatim}
    \chapter{titolo del capitolo}
    \input{file_contenente_il_testo_del_capitolo.tex}
\end{verbatim}

così da gestire gestire separatamente i vari capitoli e non avere un file unico troppo lungo da gestire.

Tutte le immagini devono essere messe nella cartella \texttt{images} (o in quella indicata come \texttt{graphicspath} nel file \texttt{main.tex}) e poi possono essere inserite comodamente nel testo con il comando: 
\begin{verbatim}
    \includegraphics[width=.5\textwidth]{nome_file.png}
\end{verbatim}

Per le citazioni è sufficiente fare \cite{felix20}, inserendo come argomento del comando \verb|\cite| lo stesso nome della citazione che è stato usato nel file \texttt{bib}. Se non si inserisce alcuna citaziona nel testo col comando \verb|\cite| non verrà generata alcuna bibliografia.

    \printbibliography

\end{document}
